\begin{teilaufgaben}
\item
The domain $\Omega$ is the square $[0,4] \times [0, 4]$.
The function $u(x,y)$ is defined on $\Omega$. 

It satisfies in $\Omega$ 
\[
\Delta u(x,y) = 10.
\]
On the boundary of $\Omega$ one has $u(x,y) = 1$. 

Determine an approximate value for $u(1,1)$ and $u(3,3)$. 

Use finite volumes \`a la Voronoi.  

\item
The domain $\Omega$ is a triangle with vertices
\[
(4,0), (0,4), (-4, 0).
\]
The function $u(x,y)$ is defined on $\Omega$.
It satisfies in $\Omega$ 
\[
\Delta u(x,y) = 0.
\]
On the boundary of $\Omega$ one has $u(x,y) = 1$ in the first quadrant
and $u(x,y) = 0$ in the second quadrant.

Determine approximate values for $u(-1,1)$ and $u(1,1)$.

Use finite volumes \`a la Voronoi.  
\end{teilaufgaben}

\begin{loesung}
\begin{teilaufgaben}
\item
By the finite volume method we have
\[
4 \cdot (1 - \tilde u(1,1)) + 4 \cdot (1 - \tilde u(1,1))
+ 2 \cdot (\tilde u(3,3) - \tilde u(1,1))
=
10 \cdot 8
\]
or
\begin{equation}
-10 \cdot \tilde u(1,1) + 2 \cdot \tilde u(3,3) = 72.
\tag{\bf 1P}
\end{equation}
By symmetry we have
\begin{equation}
-10 \cdot \tilde u(3,3) + 2 \cdot \tilde u(1,1) = 72.
\tag{\bf 1P}
\end{equation}
Therefore
\begin{equation}
(\tilde u(1,1), \tilde u(3,3)) = (-9, -9).
\tag{\bf 1P}
\end{equation}

\item
By the finite volume method we have
\[
4 \cdot (0 - \tilde u(-1,1)) + 4 \cdot (0 - \tilde u(-1,1))
+ 2 \cdot (\tilde u(1,1) - \tilde u(-1,1))
= 0
\]
or
\begin{equation}
-10 \cdot \tilde u(-1,1) + 2 \cdot \tilde u(1,1) = 0.
\tag{\bf 1P}
\end{equation}
Moreover we have
\[
4 \cdot (1 - \tilde u(1,1)) + 4 \cdot (1 - \tilde u(1,1))
+ 2 \cdot (\tilde u(-1,1) - \tilde u(1,1))
= 0
\]
or
\begin{equation}
2 \cdot \tilde u(-1,1) - 10 \cdot \tilde u(1,1) = -8.
\tag{\bf 1P}
\end{equation}
Therefore
\begin{equation}
(\tilde u(-1,1), \tilde u(1,1)) = (1/6, 5/6).
\tag{\bf 1P}
\end{equation}
\end{teilaufgaben}
\end{loesung}

